\documentclass[11pt,a4paper]{article}
\usepackage[T1]{fontenc}
\usepackage[utf8]{inputenc}
\usepackage{amsmath,amssymb,amsthm}
\usepackage[margin=1in]{geometry}
\usepackage[colorlinks=true,linkcolor=blue,urlcolor=blue]{hyperref}
\theoremstyle{plain}
\newtheorem{theorem}{Theorem}
\newtheorem{lemma}[theorem]{Lemma}
\newtheorem{proposition}[theorem]{Proposition}
\newtheorem{corollary}[theorem]{Corollary}
\theoremstyle{definition}
\newtheorem{remark}[theorem]{Remark}

\title{Recovering the unwritten recurrences of the table OEIS A238311}
\author{Adrian Perez Fontelles\\ \small Independent researcher}
\date{2 September 2026}
\begin{document}
\maketitle

\begin{abstract}
OEIS A238311 is a two-dimensional table $T(n,k)$ whose empirical recurrences are, for several of
its lines, given only as an ORDER --- ``k=2: [order 76]'' and the like --- with the
coefficients in a linked file rather than in the entry. Those conjectures can still be settled,
and the recurrences recovered. A line of the table is a fixed-width or fixed-height array
count, hence a walk count on finitely many vertices, hence satisfies some monic recurrence of
order at most that number; Berlekamp--Massey on twice as many exact terms returns the MINIMAL
one. Where its order equals the order the entry states --- and where the same holds after the
first few terms are dropped --- any recurrence of that order the line satisfies has a
characteristic polynomial that is a multiple of the minimal one and of equal degree, hence is
that polynomial. This note settles column 2.
\end{abstract}

\noindent\small 2020 Mathematics Subject Classification. 05A15, 11B37, 15A18.\normalsize

\section{The table and the unwritten conjectures}

OEIS A238311 is ``T(n,k)=Number of (n+1)X(k+1) 0..3 arrays with the maximum minus the lower median of every 2X2 subblock differing from its horizontal and vertical neighbors by exactly one.''. It is read by antidiagonals and begins
\[
256, 1776, 1776, 12372, 20128, 12372, 85984, 227622,\ \dots
\]
The entry carries, among its empirical claims:
\begin{quote}\small
k=2: [order 76]
\end{quote}
As of the ``Last modified'' line on the live entry (Jul 23 2025 10:51:22, revision 6) these are still
recorded as empirical, and the recurrences themselves are not in the entry text.

\section{Why a line of the table is a walk count}

Fix the index that the line fixes. The entry defines $T(n,k)$ as a number of arrays of a shape
determined by $n$ and $k$, subject to a condition local to a bounded window of consecutive
lines. Substituting the fixed index into the entry's own wording turns the definition into an
ordinary one-parameter array count, and for such a count the admissible windows are the
vertices of a finite digraph and an array is a walk. So the line is a walk count on $S$
vertices, and by the Cayley--Hamilton theorem it satisfies a monic integer recurrence of order
at most $S$.

\section{Recovering the recurrences}

\begin{lemma}\label{lem:bm}
A sequence known to satisfy some linear recurrence of order at most $S$ is determined, with its
minimal recurrence, by its first $2S$ terms; Berlekamp--Massey applied to them returns that
minimal recurrence exactly.
\end{lemma}

\subsection*{Column $k=2$}
Substituting the fixed index into the entry's wording gives an ordinary one-parameter array
count, a walk on $S=1776$ vertices. Berlekamp--Massey on $2S=3552$ exact terms
returns a minimal recurrence of order $76$ --- the order the entry states --- with
integer coefficients
\[
(c_1,\dots,c_{76})=(0,\ 0,\ 590,\ 6047,\ 33282,\ 143172,\ -51847,\ -3383143,\ -19098221,\ -62660439,\ 50129967,\ 1167166401,\ 3926502816,\ 4633770290,\ -30502410747,\ -161347026851,\ -92876223654,\ 780089150443,\ 1245151657046,\ 143819948233,\ -949452193524,\ -3466310431313,\ 20506201952776,\ -8028535411851,\ -375549140588226,\ -67909347802678,\ 1672222434601647,\ -1014167853197958,\ -3819901933221444,\ 12177171250980535,\ 11379879333755080,\ -19703330851331473,\ -5046728126887472,\ -107938157967819317,\ -157569757784410886,\ 466875447950918420,\ 254581489443163362,\ -971895742022921422,\ 1098815982455005785,\ 1878496926360289268,\ -4154137393506684650,\ -2489534202633012199,\ 5512403342728275355,\ 606241215183684002,\ -1593510698063116450,\ 3736090968544920876,\ -4841764212703579592,\ -6977069500653741268,\ 5774208940465613649,\ 5855695151473940574,\ -962200406677614129,\ -3635673360797615615,\ -2689068586524211097,\ 3797403507458193497,\ 2229508234425503651,\ -3424368555467162215,\ 1184789416489364860,\ 2552131844208767656,\ -1985392340164925609,\ -308502195280952876,\ 2002321901889236764,\ 344130666354013184,\ -262987165354583778,\ 271094423939197184,\ 229978471954441064,\ -55304525405563700,\ -141196591281205736,\ -16903093896289976,\ 22788803756890560,\ 11349407149777184,\ 139234319179520,\ -2069352319020288,\ -660849214850048,\ -93674083704832,\ -6774341369856,\ 34665922560).
\]
so that $a(m)=\sum_{i=1}^{76}c_i\,a(m-i)$ for every $m$. Removing the first
$76+4$ terms and repeating gives order $76$ again, which is what the
identification below needs.

\begin{theorem}
For each line above, the displayed recurrence holds for every index, and it is the recurrence
the entry records.
\end{theorem}

\begin{proof}
It holds by Lemma~\ref{lem:bm}. For the identification, the entry asserts that the line
satisfies SOME recurrence of the stated order, possibly only from some index on. Let $q$ be its
characteristic polynomial and $q_{\min}$ the minimal polynomial of the tail. Then
$q_{\min}\mid q$, and the second Berlekamp--Massey run --- on the line with its first few
terms removed --- shows $\deg q_{\min}$ equals the stated order, which is $\deg q$. Both are
monic, so $q=q_{\min}$.
\end{proof}

\section{Verification}

Three checks.

First, each line's model was compared against the entry's own published values for that line,
taken from the antidiagonal DATA read in either orientation and from the ``Table starts''
block, and agrees at every available term. This is the only point at which reading the entry's
English enters, and a misreading gives different counts.

Second, each recovered recurrence was evaluated directly on the model's values, with no
Berlekamp--Massey involved, and holds at every index tested.

Third, each order was computed twice by independent routes --- modulo a $61$-bit prime, where
every intermediate is a machine word, and again in exact rational arithmetic --- and once more
on the line with its first few terms dropped, which is what the identification requires. All
agree.

\begin{thebibliography}{9}
\bibitem{oeis} The OEIS Foundation, \emph{The On-Line Encyclopedia of Integer
Sequences}, \url{https://oeis.org}, 2026. Sequence A238311.
\bibitem{massey} J.~L.~Massey, Shift-register synthesis and BCH decoding, \emph{IEEE
Trans. Inform. Theory} \textbf{15} (1969), 122--127.
\bibitem{stanley} R.~P.~Stanley, \emph{Enumerative Combinatorics}, Volume 1, 2nd ed.,
Cambridge University Press, 2011. (Transfer-matrix method, Section 4.7.)
\bibitem{horn} R.~A.~Horn and C.~R.~Johnson, \emph{Matrix Analysis}, 2nd ed., Cambridge
University Press, 2013.
\end{thebibliography}

\end{document}
